\section{Diskussion}
\label{sec:Diskussion}
    In diesem Abschnitt sollen die Vorgehensweise und Ergebnisse der Justierung des Lasers, sowie die Funktionalität des Aufbaus kritisch betrachtet werden. 
    Während der Durchführung aufgetretene Probleme werden in Bezug zu den Ergebnissen der Auswertung gesetzt.
    \subsection{Diskussion der Laserjustierung}
        Die Einstellung der Temperatur auf die empfohlenen $\qty{50}{\degreeCelsius}$ gelang innerhalb kurzer Zeit und blieb während der Einstellung des Schwellenstroms und der vertikalen Grating-Position stabil. 
        Zwar wurde die Temperatur nur unregelmäßig kontrolliert, allerdings ließen sich zu Beginn der Messung bei Kontrollen mit kurzem Abbstand keine Schwankungen der Spannung am Multimeter feststellen.
        Die Kontrollintervalle wurden daher verlängert, da keine Abweichung zu erwarten war. Da tatsächlich keine selbstständige Veränderung der Temperatur auftrat, ist davon auszugehen, dass dieser Justierungsschritt erfolgreich war.
        \\
        Die Festlegung des Schwellenstroms für den Laserbetrieb wurde mehrmals wiederholt, um ein optimales Ergebnis zu erziehen.
        Jedoch war der Übergang zwischen LED-Betrieb und Laserbetrieb so sprunghaft, dass mit leichten Abweichungen zu rechnen ist.
        Die Messung wird deutlich beeinflusst durch die subjektive Wahrnehmung der Helligkeit des Lichtpunkts auf der IR-Sichtkarte in der Kamera. 
        Es ist zwar deutlich zu erkennen, dass der Laserbetrieb eingesetzt hat. Allerdings kann kein quantitativer Fehler bestimmt werden, der die Abweichung von der theoretischen oder tatsächlichen Laserschwelle beschreibt.
        Der Laser fiel allerdings zu keinem Zeitpunkt spontan aus dem Laser- in den LED-Betrieb.
        Daher kann hier angenommen werden, dass die gemessene Schwellenspannung von $U = \qty{3.60(0.01)}{\volt}$ nah an und mindestens über dem tatsächlichen Schwellenwert liegt.
        \\
        